\clearpage
\phantomsection
\chapter*{CONCLUSION}
\addcontentsline{toc}{chapter}{CONCLUSION}
\label{chap:conclusion}

This thesis presented the design, implementation, and evaluation of an automated game-balancing framework for a Roguelike Deckbuilder prototype, combining a Multi-Objective Evolutionary Algorithm (NSGA-II), Reinforcement Learning (RL), and a Deterministic Heuristic Agent (DHA). The system targets the central challenge in data-driven game development: finding parameter configurations that are simultaneously balanced, engaging, and internally coherent without manual iteration.

\section*{Summary and Achievements}

All five stated objectives were met.

\textbf{Objective 1 - Game Engine}: The pure Python simulation engine implements a complete Roguelike Deckbuilder prototype with 30+ cards, 20+ relics, a 15-floor procedurally generated map, and turn-based combat with 10+ status effects. Being a standard Python package, it runs natively on any platform (Windows for development, Linux on Kaggle for GPU training) with no build steps, providing the simulation backbone that the optimization layer requires.

\textbf{Objective 2 - Heuristic AI agent}: The Deterministic Heuristic Agent (DHA) provides a fast, deterministic evaluation oracle. As demonstrated in Section 4.5, the DHA achieves higher Balance and Coherence scores when used as the optimizer's evaluator compared to RL agents, precisely because its deterministic play reduces fitness estimation variance - making it the correct choice for the optimization loop.

\textbf{Objective 3 - RL AI agents}: Four MaskablePPO AI agents with distinct reward functions (Aggressive, Defensive, Balanced, Adaptive) were trained for 14 million steps each (across three curriculum phases), achieving win rates from 24.5\% to 34.0\%. Each agent occupies a distinct region of the behavioral space as confirmed by the radar chart in Section 4.2, validating that reward shaping successfully produces differentiated play styles. The Adaptive agent's highest win rate (34.0\%) demonstrates that emergent strategy switching outperforms any single explicitly designed strategy.

\textbf{Objective 4 - Game parameter optimization}: The single-objective GA improved Balance Score from 0.754 to 0.884, Engagement from 0.655 to 0.724, and Coherence from 0.853 to 0.910 in 73 minutes. NSGA-II surpassed it on every metric simultaneously - reaching Balance $= 0.960$, Engagement $= 0.900$, Coherence $= 0.984$, and increasing viable cards from 10 to 14 - in only 49 minutes (33\% faster). The best-compromise NSGA-II solution brings the average AI agent win rate to 43.5\%, within the target range of 35-45\%.

\textbf{Objective 5 - Visualization}: All evolution curves, Pareto front visualizations, and card viability analyses presented in Chapter 4 were generated by the automated visualization pipeline, confirming that the system produces interpretable, designer-readable outputs.

The primary scientific contribution is the hierarchical genome encoding, which groups approximately 500 raw game parameters into 38 semantically coherent continuous dimensions. This representation enables faster convergence, more interpretable mutations, and is generalizable to any JSON-parameterized game engine. Additional contributions include the multi-behavioral RL ensemble as a richer quality evaluator, the three-phase curriculum training strategy (easy maps $\to$ full maps $\to$ domain randomization), and the demonstration that a fully self-contained Python implementation can serve as both an ML training environment and a high-fidelity balance evaluation platform without any language bridge or external build toolchain.

\section*{Limitations}

This work has four principal limitations that bound the scope of the conclusions.

First, all evaluation relies on computational proxies - AI agents playing the game - rather than real human players. The AI agents, despite their behavioral diversity, are trained on the same optimization search space and may not represent casual players who make suboptimal decisions, or experts who discover edge-case strategies. A true player model would require human gameplay data or explicit cognitive bias modeling.

Second, the parameter hierarchy used by the HierarchicalGenome was constructed manually, based on designer intuition about which parameters are semantically related. A different hierarchy could produce different optimization results. There is no guarantee that the chosen grouping is optimal, or even locally optimal.

Third, the game prototype is smaller in scope than commercial Roguelike Deckbuilder titles: 30+ cards versus 300+ in Slay the Spire, one character class versus four, 15 floors versus 45. The system architecture is designed to scale, but the specific parameter values, genome sizes, and fitness function weights reported in Chapter 4 may not transfer directly to a larger game without re-tuning.

Fourth, the 30-generation, 50-individual optimization budget was chosen to keep experiments tractable within available compute (under 90 minutes per run on an 8-core CPU). Longer runs with larger populations would likely improve Pareto front coverage, particularly in the high-engagement region.

\section*{Future Work}

Three directions are most promising for extending this framework.

First, integrating the RL AI agents directly into the optimization loop. The current system uses the DHA as the optimizer's evaluator because RL agents are 200 times slower. However, a hybrid approach - using the DHA for initial population screening and switching to RL agents for final Pareto front refinement - would combine the speed of the DHA with the behavioral richness of the RL ensemble, potentially improving the quality of the final balance signal.

Second, automating hierarchy discovery. The hierarchical genome currently depends on a hand-crafted grouping of parameters. A meta-learning approach could analyze parameter correlation data from many GA runs to discover which parameters tend to move together, constructing the hierarchy from data rather than designer intuition. This would make the system more general and reduce the manual effort required to apply it to a new game.

Third, validating balance outputs with human playtest studies. The ultimate measure of game balance is whether human players find the game fun, fair, and engaging. Conducting structured playtests with the initial and NSGA-II-optimized configurations, measuring player-reported experience alongside AI agent metrics, would validate whether the computational proxies used in this thesis are genuinely correlated with player satisfaction.
